%%===================================================================
%% Lecture 20 -- The Quark Sector of the Standard Model:
%%               Covariant Derivative, Yukawa Diagonalization,
%%               the CKM Matrix, and the Boson Spectrum
%% 77609 -- Introduction to Elementary Particles
%% Date: 21/06/2026
%% Lecturer: Prof. Yonit Hochberg
%%===================================================================

\section{Lecture 20 -- The Quark Sector: Covariant Derivative, Yukawa Diagonalization, and the CKM Matrix}\label{sec:lec20}
\begin{center}
\begin{tabular}{ll}
  \textbf{Date:}\quad 21/06/2026 & \\
\end{tabular}
\end{center}
\hrule\vspace{1.5em}

%%------------------------------------------------------------------
\subsection{Overview}\label{subsec:lec20:overview}

In Lecture~19 we assembled the full Standard Model (SM) gauge group
$G_{\SM} = SU(3)_C\times SU(2)_L\times U(1)_Y$ and wrote its Lagrangian
schematically. Today we work through the \emph{quark sector} in detail
and discover the one structurally new feature that quarks bring and
leptons do not: a \emph{mismatch between two different mass bases} that
cannot be rotated away. This mismatch is the
\textbf{Cabibbo--Kobayashi--Maskawa (CKM) matrix}, the central result of
the lecture.

The logical chain is:
\begin{enumerate}
  \item \textbf{Field content and covariant derivative} --- recall the
        five fermion representations and build the single object,
        $D_\mu$, from which every gauge interaction is read off. We see
        how $D_\mu$ specializes to each representation.
  \item \textbf{Kinetic terms} for gauge fields, fermions, and the
        scalar.
  \item \textbf{No bare fermion masses} --- gauge invariance forbids
        $m\bar\psi_L\psi_R$ for every SM fermion.
  \item \textbf{Yukawa couplings} --- the only gauge-invariant way to
        give fermions mass, using the Higgs doublet.
  \item \textbf{Diagonalizing the Yukawa matrices} --- bi-unitary
        rotations, and the crucial observation that the up-type and
        down-type rotations of the \emph{same} left-handed doublet
        $Q_L$ differ. Their relative rotation is the CKM matrix.
  \item \textbf{Quark electric charges} from $Q = T^3 + Y$.
  \item \textbf{The boson spectrum} --- which gauge bosons get a mass
        and which stay massless.
\end{enumerate}

Throughout, ``read off from $D_\mu$'' is the operating principle:
once the representation of a field is fixed, all of its gauge couplings
are determined, not chosen.

%%==================================================================
\subsection{Field Content of the Standard Model}\label{subsec:lec20:content}

\subsubsection*{Conceptual Background}

A gauge theory is \emph{defined} by three inputs: (i) a gauge group,
(ii) the representations in which the matter fields sit, and (iii) the
pattern of spontaneous symmetry breaking. For the SM these are
\begin{equation}
  \boxed{
    \underbrace{SU(3)_C\times SU(2)_L\times U(1)_Y}_{G_{\SM}}
    \;\xrightarrow{\ \langle\phi\rangle\neq0\ }\;
    SU(3)_C\times U(1)_{\rm EM}
  }
  \label{eq:lec20:ssb}
\end{equation}
where the breaking is triggered by the Higgs field $\phi$ acquiring a
vacuum expectation value (vev). Here $C$ = color, $L$ = left
(the group acts only on left-handed fermions), and $Y$ = hypercharge.

\textit{Significance:} Everything that follows --- masses, couplings,
mixing --- is a consequence of this one line. The physics is in the
choice of representations.

The matter content is \emph{three generations} of fermions, each in
five irreducible representations, plus one Higgs scalar. The
generation index $i = 1,2,3$ labels identical copies that differ only
through their Yukawa couplings (introduced below). We write
$(\,SU(3)_C,\ SU(2)_L\,)_Y$ for the representation, giving the dimension
of each non-Abelian factor and the $U(1)_Y$ charge as a subscript:

\begin{center}
\renewcommand{\arraystretch}{1.4}
\begin{tabular}{@{}lllc@{}}
  \toprule
  Field & Name & $(\,SU(3)_C,\,SU(2)_L\,)_Y$ & Components ($i=1$)\\
  \midrule
  $Q_{Li}$ & LH quark doublet      & $(\mathbf3,\,\mathbf2)_{+1/6}$ & $(u_L,\,d_L)^T$\\
  $u_{Ri}$ & RH up-type quark      & $(\mathbf3,\,\mathbf1)_{+2/3}$ & $u_R$\\
  $d_{Ri}$ & RH down-type quark    & $(\mathbf3,\,\mathbf1)_{-1/3}$ & $d_R$\\
  $L_{Li}$ & LH lepton doublet     & $(\mathbf1,\,\mathbf2)_{-1/2}$ & $(\nu_{L},\,\ell_L)^T$\\
  $e_{Ri}$ & RH charged lepton     & $(\mathbf1,\,\mathbf1)_{-1}$   & $e_R$\\
  \midrule
  $\phi$   & Higgs doublet         & $(\mathbf1,\,\mathbf2)_{+1/2}$ & $(\phi^+,\,\phi^0)^T$\\
  \bottomrule
\end{tabular}
\end{center}

\nt{Pitfall: the \emph{left-handed} quarks live in a single
$SU(2)_L$ doublet $Q_{Li}$, but the \emph{right-handed} up and down
quarks are \emph{separate} $SU(2)_L$ singlets $u_{Ri}$ and $d_{Ri}$
with \emph{different} hypercharges ($+\tfrac23$ vs.\ $-\tfrac13$).
This left/right asymmetry is exactly what forbids bare masses
(Sec.~\ref{subsec:lec20:nomass}) and what later forces the CKM matrix
to appear (Sec.~\ref{subsec:lec20:ckm}).}

%%==================================================================
\subsection{The Covariant Derivative}\label{subsec:lec20:covd}

\subsubsection*{Mathematical Setup}

Local gauge invariance requires that ordinary derivatives
$\partial_\mu$ acting on a charged field be promoted to
\emph{covariant} derivatives $D_\mu$ that include one gauge field per
generator. With three gauge couplings $g_s$ (for $SU(3)_C$), $g$
(for $SU(2)_L$), and $g'$ (for $U(1)_Y$), the general covariant
derivative is
\begin{equation}
  \boxed{
    D_\mu = \partial_\mu
            + i g_s\, G^a_\mu\, L^a
            + i g\, W^b_\mu\, T^b
            + i g'\, B_\mu\, Y
  }
  \label{eq:lec20:covd}
\end{equation}
where:
\begin{itemize}
  \item $G^a_\mu$ ($a=1,\dots,8$) are the eight gluon fields and $L^a$
        the $SU(3)_C$ generators;
  \item $W^b_\mu$ ($b=1,2,3$) are the three weak bosons and $T^b$ the
        $SU(2)_L$ generators;
  \item $B_\mu$ is the single hypercharge boson and $Y$ the hypercharge
        (a number for each field).
\end{itemize}
The index ranges $a=1,\dots,8$ and $b=1,2,3$ are set by the dimension
of the \emph{adjoint} representation of each group, i.e.\ by the number
of generators ($8$ for $SU(3)$, $3$ for $SU(2)$, $1$ for $U(1)$).

\paragraph{How $D_\mu$ specializes to each representation.}
The objects $L^a$, $T^b$, $Y$ are abstract generators; their concrete
form depends on the representation the field sits in. The rule is:

\dfn[dfn:lec20:generators]{Generators in Each Representation}{
\begin{itemize}
  \item \textbf{$SU(3)_C$ generators $L^a$:}
        \begin{itemize}
          \item on a color \emph{triplet} (every quark field $Q_L,u_R,d_R$):
                $L^a = \dfrac{\lambda^a}{2}$, with $\lambda^a$ the
                Gell-Mann matrices;
          \item on a color \emph{singlet} (Higgs, all leptons):
                $L^a = 0$ --- they do not feel the strong force.
        \end{itemize}
  \item \textbf{$SU(2)_L$ generators $T^b$:}
        \begin{itemize}
          \item on a weak \emph{doublet} ($Q_L,L_L,\phi$):
                $T^b = \dfrac{\sigma^b}{2}$, with $\sigma^b$ the Pauli
                matrices;
          \item on a weak \emph{singlet} ($u_R,d_R,e_R$): $T^b = 0$.
        \end{itemize}
  \item \textbf{$U(1)_Y$ generator $Y$:} simply the hypercharge number
        of the field (it is always non-trivial unless $Y=0$).
\end{itemize}
}

\textit{Significance:} A field's couplings are not free inputs. Once we
declare ``$Q_L$ is a $(\mathbf3,\mathbf2)_{1/6}$,'' the explicit $D_\mu$
acting on it --- and hence all its gauge interactions --- is fully fixed.

\ex[ex:lec20:covd-QL]{Worked Example: $D_\mu$ Acting on $Q_L$}{
Take the left-handed quark doublet $Q_L\in(\mathbf3,\mathbf2)_{1/6}$.
It is a color triplet, a weak doublet, and has $Y=\tfrac16$, so each
piece of Eq.~\eqref{eq:lec20:covd} survives:
\begin{equation*}
  D_\mu Q_L
  = \left(\partial_\mu
    + i g_s\, G^a_\mu\,\frac{\lambda^a}{2}
    + i g\, W^b_\mu\,\frac{\sigma^b}{2}
    + i g'\, B_\mu\,\frac{1}{6}\right) Q_L .
\end{equation*}
By contrast, for the right-handed down quark $d_R\in(\mathbf3,\mathbf1)_{-1/3}$
the weak piece drops out (singlet $\Rightarrow T^b=0$):
\begin{equation*}
  D_\mu d_R
  = \left(\partial_\mu
    + i g_s\, G^a_\mu\,\frac{\lambda^a}{2}
    + i g'\, B_\mu\Bigl(-\tfrac13\Bigr)\right) d_R ,
\end{equation*}
and for a lepton $L_L\in(\mathbf1,\mathbf2)_{-1/2}$ the color piece drops
out ($L^a=0$). \emph{Reading the representation $\to$ writing $D_\mu$ is
the whole skill.}
}

\paragraph{The adjoint representation.}
For a field sitting in the \emph{adjoint} of a non-Abelian group (this
is where the gauge bosons themselves live), the generators are built
directly from the structure constants:
\begin{equation}
  \bigl(L^a_{\rm adj}\bigr)_{bc} = -i f^{abc}\quad(SU(3)),
  \qquad
  \bigl(T^a_{\rm adj}\bigr)_{bc} = -i\,\epsilon^{abc}\quad(SU(2)),
  \label{eq:lec20:adjoint}
\end{equation}
where $f^{abc}$ are the $SU(3)$ structure constants and
$\epsilon^{abc}$ the $SU(2)$ ones. This is the explicit statement that
\emph{in the adjoint, the generators are nothing but the structure
constants of the group} --- it is why gluons carry color and why the
$W$'s carry weak isospin (and hence self-interact, as seen in
Lecture~19).

%%==================================================================
\subsection{Kinetic Terms}\label{subsec:lec20:kinetic}

\subsubsection*{Mathematical Setup}

We organize the SM Lagrangian into four pieces,
\begin{equation}
  \lagr_{\SM} = \lagr_{\rm kin} + \lagr_{\rm Yukawa} + \lagr_{\rm scalar},
  \label{eq:lec20:lag-pieces}
\end{equation}
(with the fermion kinetic term often counted separately). The kinetic
part collects every term with derivatives. Using the field strengths
written in Lecture~19,
\begin{align}
  G^a_{\mu\nu} &= \partial_\mu G^a_\nu - \partial_\nu G^a_\mu
                  - g_s f^{abc}G^b_\mu G^c_\nu, \\
  W^a_{\mu\nu} &= \partial_\mu W^a_\nu - \partial_\nu W^a_\mu
                  - g\,\epsilon^{abc}W^b_\mu W^c_\nu, \\
  B_{\mu\nu}   &= \partial_\mu B_\nu - \partial_\nu B_\mu,
\end{align}
the kinetic Lagrangian is
\begin{equation}
  \boxed{
  \begin{aligned}
    \lagr_{\rm kin}
    &= -\tfrac14 G^a_{\mu\nu}G^{a\mu\nu}
       -\tfrac14 W^a_{\mu\nu}W^{a\mu\nu}
       -\tfrac14 B_{\mu\nu}B^{\mu\nu}\\
    &\quad + \sum_{\psi} i\,\bar\psi\,\slashed{D}\,\psi
           + (D_\mu\phi)^\dagger(D^\mu\phi),
  \end{aligned}}
  \label{eq:lec20:kinetic}
\end{equation}
where the sum runs over \emph{all} fermion fields,
\begin{equation}
  \psi \in \{\,Q_{Li},\ u_{Ri},\ d_{Ri},\ L_{Li},\ e_{Ri}\,\},
  \qquad i=1,2,3,
  \label{eq:lec20:psi-sum}
\end{equation}
and for each field $D_\mu$ takes the representation-specific form of
Sec.~\ref{subsec:lec20:covd}. The Feynman-slash notation is
$\slashed{D}=\gamma^\mu D_\mu$.

\textit{Significance:} A single compact line,
$\sum_\psi i\bar\psi\slashed D\psi$, encodes \emph{every} gauge
interaction of \emph{every} fermion --- strong, weak, and
electromagnetic --- because $D_\mu$ knows each field's representation.

%%==================================================================
\subsection{No Bare Mass Terms for Fermions}\label{subsec:lec20:nomass}

\subsubsection*{Conceptual Background}

A Dirac mass term couples a left-handed field to a right-handed one,
\begin{equation}
  \lagr_{\rm mass} = -m\,\bar\psi_L\psi_R + \text{h.c.}
  \label{eq:lec20:diracmass}
\end{equation}
For this to be allowed in the Lagrangian it must be a \emph{singlet}
under the full gauge group --- otherwise it would not be gauge
invariant. The question is therefore purely group-theoretic: can
$\bar\psi_L\psi_R$ form a singlet for any SM fermion?

\clm[clm:lec20:nomass]{Bare Fermion Masses Are Forbidden}{
For every SM fermion, the product $\bar\psi_L\psi_R$ fails to be a
gauge singlet, so
\begin{equation}
  \boxed{m_{\rm bare} = 0 \quad\text{for all SM fermions.}}
  \label{eq:lec20:nomass-result}
\end{equation}
}

\subsubsection*{Why it fails --- the quark case}

Try to build a quark mass term, e.g.\ $\bar Q_L\,u_R$:
\begin{itemize}
  \item \textbf{$SU(3)_C$:} $\bar Q_L\sim\bar{\mathbf3}$,
        $u_R\sim\mathbf3$, and $\bar{\mathbf3}\otimes\mathbf3\supset\mathbf1$.
        \checkmark\ (color is fine).
  \item \textbf{$SU(2)_L$:} $\bar Q_L\sim\bar{\mathbf2}$, $u_R\sim\mathbf1$,
        and $\bar{\mathbf2}\otimes\mathbf1 = \mathbf2\neq\mathbf1$.
        ($\boldsymbol\times$) (no singlet --- this term has a leftover weak index).
  \item \textbf{$U(1)_Y$:} $Y(\bar Q_L)+Y(u_R) = -\tfrac16+\tfrac23
        = +\tfrac12\neq 0$. ($\boldsymbol\times$) (hypercharge does not cancel).
\end{itemize}
Two of the three factors fail. The same happens for the leptons:
$\bar L_L e_R$ is an $SU(2)_L$ doublet ($\mathbf2$, not a singlet) and
carries net hypercharge $-\tfrac12+(-1)\neq0$.

\textit{Significance:} Chirality plus the gauge assignment makes a bare
mass impossible. \emph{All} fermion masses must therefore come from
somewhere else --- the Higgs mechanism via Yukawa couplings. This is
not an aesthetic choice; it is forced.

\nt{Pitfall: it is the \emph{mismatch} between the left-handed
(doublet) and right-handed (singlet) assignments that kills the mass
term. If left and right fields lived in the \emph{same} representation
(as in QED or QCD alone, which are vector-like), a bare mass would be
allowed. The chiral structure of the weak interaction is exactly what
makes the Higgs necessary.}

%%==================================================================
\subsection{Yukawa Couplings}\label{subsec:lec20:yukawa}

\subsubsection*{Physical Motivation}

Since $\bar\psi_L\psi_R$ alone is not a singlet, we repair it by
inserting the Higgs doublet $\phi$ to soak up the leftover $SU(2)_L$
and $U(1)_Y$ charges. The result is a \emph{Yukawa coupling}: a
gauge-invariant, dimension-4 (renormalizable) term linking two
fermions to the scalar. When $\phi$ acquires its vev, each Yukawa term
becomes a mass term.

\subsubsection*{Mathematical Setup}

We need both $\phi$ and a partner that carries the opposite weak
charge. Define the \emph{charge-conjugate} Higgs doublet
\begin{equation}
  \tilde\phi \equiv i\sigma^2\phi^* ,
  \qquad (\tilde\phi)_a = \epsilon_{ab}\,\phi_b^* ,
  \label{eq:lec20:phitilde}
\end{equation}
which is also an $SU(2)_L$ doublet but with hypercharge $Y(\tilde\phi)
= -\tfrac12$ (opposite to $\phi$). With $\phi$ and $\tilde\phi$
available we can write all three Yukawa terms:
\begin{equation}
  \boxed{
    -\lagr_{\rm Yukawa}
    = Y^u_{ij}\,\bar Q_{Li}\,\tilde\phi\,u_{Rj}
    + Y^d_{ij}\,\bar Q_{Li}\,\phi\,d_{Rj}
    + Y^e_{ij}\,\bar L_{Li}\,\phi\,e_{Rj}
    + \text{h.c.}
  }
  \label{eq:lec20:yukawa}
\end{equation}
where $Y^u, Y^d, Y^e$ are \emph{arbitrary complex $3\times3$ matrices}
in generation space ($i,j=1,2,3$).

\paragraph{Checking gauge invariance through hypercharge.}
The quickest consistency check is $U(1)_Y$: the hypercharges of the
three factors must sum to zero.
\begin{align*}
  \text{up:}   \quad & Y(\bar Q_L) + Y(\tilde\phi) + Y(u_R)
                       = -\tfrac16 - \tfrac12 + \tfrac23 = 0,\quad\checkmark\\
  \text{down:} \quad & Y(\bar Q_L) + Y(\phi) + Y(d_R)
                       = -\tfrac16 + \tfrac12 - \tfrac13 = 0,\quad\checkmark\\
  \text{lepton:}\quad& Y(\bar L_L) + Y(\phi) + Y(e_R)
                       = +\tfrac12 + \tfrac12 - 1 = 0.\quad\checkmark
\end{align*}
The up-type term \emph{must} use $\tilde\phi$ (not $\phi$) precisely
because only $\tilde\phi$ carries the hypercharge that makes the sum
vanish.

\textit{Significance:} The Yukawa couplings are the \emph{only}
renormalizable terms that can give fermions mass, and they are
unavoidably matrices in generation space. Their off-diagonal entries
are the seed of all flavor physics.

\nt{Pitfall: do not confuse the conjugate doublet $\tilde\phi$ with
$\phi^\dagger$. $\tilde\phi = i\sigma^2\phi^*$ is still a
\emph{column} doublet (it transforms as a $\mathbf2$ of $SU(2)_L$),
whereas $\phi^\dagger$ is a row. Only $\tilde\phi$ can multiply $u_R$
to make an invariant.}

\paragraph{The lepton sector is the easy case.}
The leptons are exactly as in the Lepton Standard Model of
Lectures~17--18, with the only change that they are now also $SU(3)_C$
singlets (which does nothing). As before, a single bi-unitary rotation
puts $Y^e$ into real diagonal form,
\begin{equation}
  Y^e \;\xrightarrow{\ \text{rotate}\ }\;
  \hat Y^e = \mathrm{diag}(y_e,\,y_\mu,\,y_\tau),
  \label{eq:lec20:ye-diag}
\end{equation}
and in this \emph{mass basis} the doublet components are named
$(\nu_{eL},e_L),(\nu_{\mu L},\mu_L),(\nu_{\tau L},\tau_L)$. There is no
mixing among leptons. The quark sector is where something genuinely new
happens.

%%==================================================================
\subsection{Diagonalizing the Quark Yukawa Matrices}\label{subsec:lec20:ckm}

\subsubsection*{Mathematical Setup: Bi-Unitary Diagonalization}

Any complex matrix $M$ can be brought to real, non-negative diagonal
form by a \emph{bi-unitary} transformation: there exist unitary
matrices $V_L, V_R$ such that $V_L M V_R^\dagger$ is diagonal with
non-negative entries (this is the singular-value decomposition). We
apply this independently to $Y^u$ and $Y^d$.

\paragraph{Step 1 --- diagonalize the up-type Yukawa.}
Choose unitary matrices $V_{uL}, V_{uR}$ so that
\begin{equation}
  \hat Y^u = V_{uL}\,Y^u\,V_{uR}^\dagger
           = \mathrm{diag}(y_u,\,y_c,\,y_t),
  \label{eq:lec20:yu-diag}
\end{equation}
real and diagonal, ordered by size $y_u<y_c<y_t$. The rotation
$V_{uL}$ acts on the \emph{up-components} of the left-handed doublets
$Q_L$, and $V_{uR}$ on the right-handed fields $u_R$. In this basis we
name the resulting mass-eigenstate fields
\begin{equation*}
  \text{LH up-components:}\ (u_L,\,c_L,\,t_L),
  \qquad
  \text{RH up quarks:}\ (u_R,\,c_R,\,t_R),
\end{equation*}
ordered to match the eigenvalues (up, charm, top).

\paragraph{Step 2 --- diagonalize the down-type Yukawa.}
\emph{Alternatively}, choose a \emph{different} pair $V_{dL}, V_{dR}$ so
that
\begin{equation}
  \hat Y^d = V_{dL}\,Y^d\,V_{dR}^\dagger
           = \mathrm{diag}(y_d,\,y_s,\,y_b),
  \label{eq:lec20:yd-diag}
\end{equation}
real and diagonal, ordered $y_d<y_s<y_b$ (down, strange, bottom). Here
$V_{dL}$ acts on the \emph{down-components} of $Q_L$ and $V_{dR}$ on the
$d_R$ fields.

\subsubsection*{The Key Tension: One Doublet, Two Rotations}

Here is the crux. The \emph{same} left-handed doublet $Q_{Li}$ supplies
\emph{both} the up-components (which appear in $Y^u$) and the
down-components (which appear in $Y^d$). To diagonalize $Y^u$ we must
rotate $Q_L$ by $V_{uL}$; to diagonalize $Y^d$ we must rotate the
\emph{same} $Q_L$ by $V_{dL}$. In general these are different unitary
matrices:
\begin{equation}
  \boxed{V_{uL}\neq V_{dL}.}
  \label{eq:lec20:vneq}
\end{equation}
We therefore \emph{cannot} make both $Y^u$ and $Y^d$ diagonal at the
same time. Concretely:
\begin{itemize}
  \item In the basis where $Y^u$ is diagonal (the up-quark mass basis),
        the down Yukawa is \emph{not} diagonal: it equals its own
        diagonal form rotated by a leftover matrix $V$,
        $\;Y^d \to \hat Y^d\,V$.
  \item In the basis where $Y^d$ is diagonal (the down-quark mass
        basis), $Y^u$ is rotated by the same $V$.
\end{itemize}
This leftover rotation is the relative misalignment of the two
diagonalizations,
\begin{equation}
  \boxed{V_{\rm CKM} = V_{uL}\,V_{dL}^\dagger.}
  \label{eq:lec20:ckm-def}
\end{equation}

\dfn[dfn:lec20:ckm]{Cabibbo--Kobayashi--Maskawa (CKM) Matrix}{
The CKM matrix $V_{\rm CKM}=V_{uL}V_{dL}^\dagger$ is the unitary
$3\times3$ matrix measuring the mismatch between the up-quark and
down-quark mass bases. It is named for \textbf{C}abibbo (who found the
two-generation version), \textbf{K}obayashi, and \textbf{M}askawa (who
extended it to three generations and showed it can accommodate CP
violation). It is a \emph{physical}, measurable object and plays a
\emph{crucial role in charged-current weak interactions} (the $W^\pm$
couplings to quarks).
}

\textit{Significance:} The mass basis for up-type quarks is
\emph{not} the mass basis for down-type quarks. Their unavoidable
relative rotation $V_{\rm CKM}$ is the origin of quark flavor mixing
and of all CP violation in the SM. Leptons have no analogue here
because (with massless neutrinos) only one charged-lepton Yukawa needs
diagonalizing.

\subsubsection*{Why this matters physically (foreshadowing Thursday)}

The CKM matrix has no effect on the photon, $Z$, or gluon couplings:
those are \emph{flavor-diagonal} neutral currents, and a unitary
rotation acting identically on a fermion and its conjugate leaves them
unchanged (no tree-level flavor-changing neutral currents). It survives
only in the \emph{charged current}, where a $W^\pm$ connects an
up-type quark to a down-type quark of the \emph{other} mass basis:
\begin{equation}
  \lagr_{\rm CC}
  \supset
  -\frac{g}{\sqrt2}\,W^+_\mu\,
  \bar u_{Li}\,\gamma^\mu\,(V_{\rm CKM})_{ij}\,d_{Lj}
  + \text{h.c.}
  \label{eq:lec20:cc-foreshadow}
\end{equation}
So if a $W$ decay leaves you holding one flavor of quark, there is an
amplitude --- weighted by the corresponding CKM element --- to find a
\emph{different} down-type flavor as its partner. This is how flavor
``changes'' even though each individual vertex is diagonal in the weak
basis. (We will derive Eq.~\eqref{eq:lec20:cc-foreshadow} carefully in
the next lecture.)

\subsubsection*{Method: Finding the Physical Mixing Matrix}
\begin{enumerate}
  \item Write the general complex Yukawa matrices $Y^u, Y^d$ in the weak
        (interaction) basis.
  \item Bi-unitarily diagonalize each:
        $\hat Y^u = V_{uL}Y^u V_{uR}^\dagger$,
        $\hat Y^d = V_{dL}Y^d V_{dR}^\dagger$.
  \item Identify the four left/right rotations
        $V_{uL},V_{uR},V_{dL},V_{dR}$.
  \item Form the physical mixing matrix from the two \emph{left-handed}
        rotations that act on the shared doublet:
        $V_{\rm CKM} = V_{uL}V_{dL}^\dagger$.
  \item The right-handed rotations $V_{uR},V_{dR}$ are unphysical: they
        act on singlets that never share an $SU(2)_L$ structure, and
        drop out of all observables.
\end{enumerate}

\ex[ex:lec20:cabibbo]{Worked Example: Two-Generation (Cabibbo) Mixing}{
With only two generations, the leftover misalignment is a single
real rotation angle, the \emph{Cabibbo angle} $\theta_c$:
\begin{equation*}
  V_{\rm CKM}^{(2)}
  = \begin{pmatrix} \cos\theta_c & \sin\theta_c\\
                    -\sin\theta_c & \cos\theta_c \end{pmatrix},
  \qquad \theta_c\approx 13^\circ .
\end{equation*}
The charged current then reads, schematically,
\begin{equation*}
  \bar u_L\gamma^\mu(\cos\theta_c\, d_L + \sin\theta_c\, s_L)
  + \bar c_L\gamma^\mu(-\sin\theta_c\, d_L + \cos\theta_c\, s_L).
\end{equation*}
\textbf{Interpretation:} the up quark couples mostly to the down quark
($\cos\theta_c\approx0.97$) but with a small amplitude
($\sin\theta_c\approx0.22$) to the strange quark. This single number
explains why strangeness-changing decays (e.g.\
$K^+\to\mu^+\nu_\mu$) are suppressed relative to
strangeness-conserving ones --- they pay a factor $\sin^2\theta_c$ in
rate. Counting the parameters: a general $2\times2$ unitary matrix has
4 real parameters, but 3 of its phases can be absorbed by redefining
the quark fields, leaving exactly \emph{one} physical angle.
}

\nt{Pitfall: the right-handed rotations $V_{uR}$ and $V_{dR}$ are
\emph{not} physical and do \emph{not} enter $V_{\rm CKM}$. A common
error is to write $V_{\rm CKM}=V_{uL}V_{dR}^\dagger$ or to include a
right-handed mixing. Only the left-handed doublet is shared between the
two Yukawa terms, so only $V_{uL}$ and $V_{dL}$ can fail to cancel.}

\nt{Pitfall: parameter counting for three generations. A unitary
$3\times3$ matrix has $9$ real parameters. Of these, $5$ relative quark
phases can be rephased away, leaving $4$ physical parameters:
\emph{three mixing angles and one CP-violating phase}. That surviving
phase --- impossible with only two generations --- is Kobayashi and
Maskawa's reason for predicting a third generation.}

%%==================================================================
\subsection{Electric Charges of the Quarks}\label{subsec:lec20:charges}

\subsubsection*{Conceptual Background}

After spontaneous symmetry breaking, the unbroken electromagnetic
generator is the Gell-Mann--Nishijima combination
\begin{equation}
  Q = T^3 + Y,
  \label{eq:lec20:gmn}
\end{equation}
which survives because the Higgs vev sits in the lower (electrically
neutral) component of the doublet, leaving $Q\langle\phi\rangle=0$.
We now apply $Q=T^3+Y$ to each quark field, using $T^3=\pm\tfrac12$ for
the upper/lower entry of an $SU(2)_L$ doublet and $T^3=0$ for a singlet.

For the left-handed doublet $Q_L$ with $Y=\tfrac16$:
\begin{align}
  Q(u_L) &= T^3 + Y = +\tfrac12 + \tfrac16 = +\tfrac23,
  \label{eq:lec20:Qu}\\
  Q(d_L) &= T^3 + Y = -\tfrac12 + \tfrac16 = -\tfrac13.
  \label{eq:lec20:Qd}
\end{align}
For the right-handed singlets ($T^3=0$, so $Q=Y$):
\begin{equation}
  Q(u_R) = Y = +\tfrac23,
  \qquad
  Q(d_R) = Y = -\tfrac13.
  \label{eq:lec20:QuRdR}
\end{equation}

\clm[clm:lec20:charges]{Quark Electric Charges}{
The left- and right-handed up-type quarks both have
$Q=+\tfrac23$, and the down-type quarks both have $Q=-\tfrac13$:
\begin{equation}
  \boxed{
    Q(u,c,t) = +\tfrac23,
    \qquad
    Q(d,s,b) = -\tfrac13.
  }
  \label{eq:lec20:quark-charges}
\end{equation}
}

\textit{Significance:} The fractional charges $+\tfrac23$ and
$-\tfrac13$ are not inputs; they fall out of the hypercharge assignments
$Y(Q_L)=\tfrac16$, $Y(u_R)=\tfrac23$, $Y(d_R)=-\tfrac13$ through
$Q=T^3+Y$. That $u_L$ and $u_R$ end up with the \emph{same} $Q$ (and
likewise $d_L,d_R$) is exactly what makes the photon coupling
vector-like and a Dirac mass electrically consistent.

\nt{Pitfall: the up-type and down-type quarks have the same hypercharge
\emph{inside the doublet} ($Y=\tfrac16$ for both $u_L$ and $d_L$) but
differ in $T^3$; it is $T^3$ that splits their electric charges by one
unit. This unit of charge is exactly what the $W^\pm$ (charge $\pm1$)
carries away in a charged-current transition $u\leftrightarrow d$.}

%%==================================================================
\subsection{The Boson Spectrum}\label{subsec:lec20:spectrum}

\subsubsection*{Physical Motivation}

With the matter sector understood, we ask which gauge bosons are
massive and which remain massless after
$SU(2)_L\times U(1)_Y\to U(1)_{\rm EM}$ (with $SU(3)_C$ untouched).
The rule from spontaneous symmetry breaking: a gauge boson acquires a
mass for each \emph{broken} generator, by ``eating'' the corresponding
would-be Goldstone boson; gauge bosons of \emph{unbroken} generators
stay massless.

\subsubsection*{Counting}

The electroweak factor $SU(2)_L\times U(1)_Y$ has $3+1=4$ generators.
The breaking leaves one unbroken combination ($U(1)_{\rm EM}$), so
\begin{equation}
  \boxed{
    4 \text{ EW generators}
    \;\longrightarrow\;
    \underbrace{3}_{\text{broken}\ \to\ \text{massive}}
    + \underbrace{1}_{\text{unbroken}\ \to\ \text{massless}}.
  }
  \label{eq:lec20:ew-count}
\end{equation}
The three broken generators give three massive bosons --- $W^+$, $W^-$,
$Z$ --- which absorb the three would-be Goldstone bosons of the Higgs
doublet as their longitudinal polarizations. The one unbroken
generator gives the massless photon $A$ (or $A^0$). Since $SU(3)_C$ is
not broken at all, its eight gluons $G^a_\mu$ remain massless.

\begin{center}
\renewcommand{\arraystretch}{1.3}
\begin{tabular}{@{}llc@{}}
  \toprule
  Boson(s) & Origin & Mass\\
  \midrule
  $8$ gluons $G^a_\mu$ & unbroken $SU(3)_C$ adjoint & massless\\
  $W^+,\,W^-$          & 2 broken EW generators     & massive\\
  $Z$                  & 1 broken EW generator      & massive\\
  photon $A$           & unbroken $U(1)_{\rm EM}$    & massless\\
  \bottomrule
\end{tabular}
\end{center}

\subsubsection*{Goldstone bookkeeping and the Higgs}

The complex Higgs doublet $\phi=(\phi^+,\phi^0)^T$ has four real
degrees of freedom. Three are eaten:
\begin{equation}
  \underbrace{4}_{\text{real d.o.f.\ of }\phi}
  = \underbrace{3}_{\text{Goldstones eaten by }W^\pm,Z}
  + \underbrace{1}_{\text{physical Higgs }h}.
  \label{eq:lec20:goldstone-count}
\end{equation}
The three eaten Goldstones become the longitudinal components of
$W^\pm$ and $Z$; the remaining degree of freedom is the physical,
massive, real scalar Higgs boson $h$.

\subsubsection*{The $\rho$ parameter}

Because the breaking is driven by an $SU(2)_L$ \emph{doublet} (the same
as in the Lepton Standard Model), the tree-level relation between the
$W$ and $Z$ masses is preserved:
\begin{equation}
  \boxed{
    \rho \equiv \frac{m_W^2}{m_Z^2\cos^2\theta_W} = 1 .
  }
  \label{eq:lec20:rho}
\end{equation}

\textit{Significance:} Adding color ($SU(3)_C$) and the quarks does
\emph{not} disturb the electroweak boson spectrum --- that part of the
story is identical to the Lepton Standard Model. The genuinely new
ingredient relative to the LSM is the eight massless gluons, protected
by the unbroken color symmetry, and the quarks themselves, whose up-
and down-type mass bases differ (the CKM story above).

%%------------------------------------------------------------------
\subsection*{To Remember}

\begin{itemize}
  \item \textbf{Covariant derivative:}
        $D_\mu = \partial_\mu + ig_s G^a_\mu L^a + ig W^b_\mu T^b
        + ig' B_\mu Y$; specialize per representation
        ($L^a=\lambda^a/2$ on color triplets, $0$ on singlets;
        $T^b=\sigma^b/2$ on weak doublets, $0$ on singlets).
  \item \textbf{Adjoint generators:} $(T^a_{\rm adj})_{bc}=-if^{abc}$ ---
        gauge bosons carry the charge of their group and self-interact.
  \item \textbf{No bare masses:} $\bar\psi_L\psi_R$ is never a gauge
        singlet (fails $SU(2)_L$ and/or $U(1)_Y$); all masses come from
        Yukawa couplings.
  \item \textbf{Yukawa terms:}
        $Y^u_{ij}\bar Q_{Li}\tilde\phi u_{Rj}
        + Y^d_{ij}\bar Q_{Li}\phi d_{Rj}
        + Y^e_{ij}\bar L_{Li}\phi e_{Rj} + \text{h.c.}$,
        with $\tilde\phi=i\sigma^2\phi^*$ for up-type; $Y$'s are general
        complex $3\times3$.
  \item \textbf{CKM matrix:} $V_{\rm CKM}=V_{uL}V_{dL}^\dagger$ --- the
        up-quark and down-quark mass bases are misaligned because the
        shared doublet $Q_L$ needs two different left-rotations
        ($V_{uL}\neq V_{dL}$). It enters \emph{only} the charged current;
        $3$ angles $+1$ CP phase for three generations.
  \item \textbf{Quark charges:} $Q=T^3+Y\Rightarrow Q(u,c,t)=+\tfrac23$,
        $Q(d,s,b)=-\tfrac13$.
  \item \textbf{Boson spectrum:} $SU(2)_L\times U(1)_Y\to U(1)_{\rm EM}$
        gives massive $W^\pm,Z$ (3 broken generators, 3 eaten
        Goldstones) and massless $\gamma$; $SU(3)_C$ unbroken $\Rightarrow$
        $8$ massless gluons; $1$ physical Higgs $h$; $\rho=1$ at tree
        level.
\end{itemize}

%%------------------------------------------------------------------
\begin{furtherreading}
\subsection*{Further Reading}

\begin{itemize}
  \item D.\,J. Griffiths, \textit{Introduction to Elementary Particles}
        (2nd ed.): Ch.~9--10, for quark charges, the Cabibbo angle, and
        an introductory treatment of quark mixing.
  \item F. Halzen \& A.\,D. Martin, \textit{Quarks and Leptons}:
        Ch.~12, on the GIM mechanism, the Cabibbo angle, and the CKM
        matrix in charged-current processes.
  \item Y. Grossman \& Y. Nir, \textit{The Standard Model: From
        Fundamental Physics to Experimental Probes}: Ch.~9--11
        (\textit{The Standard Model with quarks}); this is the closest
        match to the notation and logic of this course, including the
        bi-unitary diagonalization and the $V_{\rm CKM}=V_{uL}V_{dL}^\dagger$
        derivation.
  \item M.\,E. Peskin \& D.\,V. Schroeder, \textit{An Introduction to
        Quantum Field Theory}: \S20.3, for the quark Yukawa sector and
        the CKM matrix.
  \item Particle Data Group, \textit{Review of Particle Physics} (2024):
        \textit{CKM Quark-Mixing Matrix} review for measured magnitudes,
        angles, and the Jarlskog invariant. \url{https://pdg.lbl.gov/}
\end{itemize}
\end{furtherreading}
